\section{API 定价：生态公共品与商业替代}

Reddit API 同时服务版主机器人、研究工具、无障碍客户端、爱好项目、第三方应用和大规模商业数据使用。统一价格看似简单，却会伤害低成本高公共价值工具；全部免费又可能让平台承担不断上升的基础设施成本，并失去广告或数据许可收入。真正的设计问题是如何分类、计量、通知和迁移，而不只是“收费是否正确”。

\subsection{把使用者按行为而非身份分类}

\term{Rate limit} 是对单位时间调用次数、并发量或数据量的限制，用于保护容量和公平性。API 设计可以给版主工具与无障碍应用更高公共价值权重，为研究和爱好项目提供免费额度，对会替代官方客户端或批量提取数据的商业使用建立合同。分类标准必须可观察，否则团队只能靠公司名气和谈判能力决定价格。

\begin{figure}[H]
\centering
\includegraphics[width=0.96\textwidth]{images/08-api-economics.png}
\caption{API 经济的分层：社区工具、研究、第三方客户端与大规模商业使用需要不同合同。}
\figsource{依据访谈字幕 32:55--38:00 重绘，`images/08-api-economics.png`。}
\end{figure}

\begin{importantbox}{读图：价格之前先定义替代关系}
Moderation bot 增强 Reddit 本身，第三方完整客户端可能替代官方广告和产品体验，批量模型训练则把语料价值迁移到外部系统。调用次数相同，经济含义可能不同。一个可辩护的制度应公开类别、免费额度、性能保证、例外申请和迁移窗口，而不是在最后一刻用单一单价处理全部使用者。
\end{importantbox}

\subsection{抗议是反馈，不是自动否决权}

API 变更引发 subreddit 关闭和社区抗议。Huffman 在访谈中承认抗议符合 Reddit 的民主属性，同时认为管理层仍需对平台可持续性作出决定。两者可以同时成立：社区有权表达集体成本，平台也必须评估安全、收入、基础设施和长期产品控制。问题在于决策过程能否让受影响者提前理解理由、测试替代方案并获得合理迁移时间。

\begin{knowledgebox}{课堂提示：不要把冲突压成善恶叙事}
第三方开发者投入多年建立产品，用户形成依赖，版主工具承担平台公共劳动；Reddit 则支付存储、带宽、安全和工程成本，并面临商业替代。冲突来自旧的隐含合同被重新定价。最有用的复盘不是判断谁更愤怒，而是找出哪些依赖没有被登记、哪些群体没有迁移路径、哪些成本从未透明。
\end{knowledgebox}

\begin{table}[H]
\centering
\begin{tabular}{L{0.2\textwidth}L{0.34\textwidth}L{0.36\textwidth}}
\toprule
阶段 & 平台应提供 & 生态应提供 \\
\midrule
发现 & 使用分类、成本模型、依赖清单 & 真实调用量、关键功能、无障碍需求 \\
设计 & 价格、免费层、SLA、例外条件 & 替代方案和可接受迁移期 \\
试点 & 沙盒、监控、回滚条件 & 小流量测试和兼容反馈 \\
迁移 & 文档、工具、支持与明确日期 & 客户端更新、用户通知、数据导出 \\
复盘 & 中断、投诉、收入和生态损失 & 失败案例、遗留依赖与改进建议 \\
\bottomrule
\end{tabular}
\caption{API 政策变更是一项平台迁移工程，不是一封定价公告。}
\end{table}

\begin{warningbox}{上市后的约束会改变决策口径}
公开公司需要说明收入、成本与风险，但股东责任不能自动覆盖社区责任。志愿版主、内容创作者和长期开发者贡献了平台价值，却不一定出现在传统成本表中。成熟管理需要把生态健康纳入可观察指标，而不是等抗议发生后才发现依赖。
\end{warningbox}

\subsection{本章小结}

API 争议的根因是不同使用模式共享了一个历史上近似免费的接口。可持续方案需要行为分类、透明计量、公共价值例外、商业合同和迁移程序。开放性不是永远零价格，可持续性也不等于突然关闭；二者之间需要长期可预测的接口治理。

